\section{Dataset construction and graph extraction}
\label{app:datasets}

Every corpus in this paper is reduced to the same intermediate object: a
\emph{typed attributed graph} in which each node is a document (an article, a
source file, or a paper) carrying its full text, and each directed edge is an
in-corpus reference recovered from that text (a wikilink, an \texttt{import}, or
a \texttt{\textbackslash cite}). The masks of \Cref{sec:method} consume this
graph, so the fidelity of edge extraction bounds what the model can learn. The
pipeline is uniform across modalities:
\[
\text{extract} \;\to\; \texttt{graph.jsonl}\;(+\text{content}) \;\to\;
\text{pretokenize} \;\to\; \texttt{tokenized\_graph.jsonl}+\texttt{shard\_*.bin}
\;\to\; \text{split}.
\]
Only extraction is modality-specific; normalization, pretokenization, and
splitting are shared. A load-bearing invariant throughout is that
\textbf{link positions are never stored}: the graph records only that an edge
exists, and the token span realizing a link is re-detected at runtime by the
online detectors (\Cref{app:linkdet}). Extraction and detection must therefore
agree on what constitutes a link.

\subsection{Per-modality edge extraction}

\paragraph{Wikipedia (and Wikisource).}
Extraction is two-stage and text-mediated. A dump extractor consumes the
CirrusSearch JSON export, keeps namespace-0 articles, and writes one Markdown
file per article, rewriting wikitext links \texttt{[[target|label]]} to Markdown
\texttt{[label](target)} with the \emph{target left as the raw article title}
(spaces preserved, no hashing). A second pass re-parses each file, extracts and
URL-unquotes link targets, normalizes them (\Cref{app:datasets:norm}), and emits
an edge \emph{iff} the normalized target is in the node set. Outgoing edges
retain \emph{all} targets (unresolved ones survive as latent cross-source
danglers), whereas incoming adjacency is restricted to in-set sources. The
fraction of wikilinks that resolve to an in-corpus node is \emph{not logged}
(\Cref{app:datasets:attack}).

\paragraph{The Stack / Python.}
Edges are \texttt{import} dependencies. A \texttt{tree-sitter} pass extracts
\texttt{import} and \texttt{from~\ldots~import} statements from each file in The
Stack \citep{kocetkov2022stack}; a hardcoded denylist (roughly three dozen names:
the standard library plus ubiquitous third-party packages such as \texttt{numpy},
\texttt{pandas}, \texttt{torch}, \texttt{django}) drops non-local modules. Each
surviving reference is resolved to a file \emph{within the same repository} by
walking dotted module paths (mapping \texttt{\_\_init\_\_} to its parent
package), preferring an exact \texttt{.py} match, then an exact module match,
then a prefix match. When several candidates remain the tie is broken by an
\emph{unseeded} \texttt{random.choice}, so import resolution is \textbf{not
deterministic across runs}. Only intra-repository edges are produced (node id
\texttt{repo:path}), and repositories are hash-partitioned into 256 buckets so a
repo is never split across shards. Files with fewer than two intra-repo links are
then discarded.

\paragraph{arXiv / unarXive.}
Edges are citations, extracted from unarXive \citep{saier2020unarxive,
saier2023unarxive}. Each \texttt{\textbackslash cite} is resolved by two
mechanisms: unarXive's own arXiv-id annotations resolve a direct minority
($\approx 14.5\%$), and an OpenAlex identifier map \citep{priem2022openalex}
lifts total coverage to $\approx 66\%$. An edge is emitted only when the cited
work resolves to an id \emph{present in the corpus}; citations that resolve
outside the corpus are \emph{deleted from the text entirely} (no dangling marker,
no edge), so the model never sees a citation it cannot follow, while in-corpus
citations are rewritten to \texttt{\textbackslash cite\{title\}}. Title
collisions are disambiguated by appending \texttt{ (arxiv:id)} and affect
$\approx 0.13\%$ of citations.

\begin{table}[t]
  \centering
  \small
  \caption{Per-corpus graph scale for the specific samples used here, after
  extraction and filtering; these are as-built figures, not intrinsic corpus
  properties. Wikipedia figures are post-resolution; the Stack figure is the
  100M-file working set after the two-intra-repo-link filter, which retains
  $28.7\%$ of files (this fraction grows with sample size---$5.3\%$/$25.1\%$/
  $28.7\%$ at the 1M/10M/100M scales).}
  \label{tab:datasets:scale}
  \begin{tabular}{llrrr}
    \toprule
    Corpus & Edge type & Nodes & Edges & Mean out-deg. \\
    \midrule
    \texttt{enwiki}       & wikilink & 275{,}199 & 2{,}321{,}731 & 8.44 \\
    \texttt{enwikisource} & wikilink & 662{,}058 &   114{,}523 & 0.17 \\
    The Stack (Python)    & import   & 3{,}560{,}799 & 6{,}335{,}322 & 1.78 \\
    arXiv / unarXive      & citation & \multicolumn{3}{c}{$\approx 66\%$ of \texttt{\textbackslash cite} resolved in-corpus} \\
    \bottomrule
  \end{tabular}
\end{table}

The enwiki graph is dense (mean out-degree $8.44$) and enwikisource very sparse
(mean $0.17$), consistent with the reference topology of WikiLinkGraphs
\citep{consonni2019wikilinkgraphs}. On the 100M Stack sample, $\approx 34\%$ of
nodes have zero in-degree and $\approx 25\%$ zero out-degree.

\subsection{Identifier normalization}
\label{app:datasets:norm}

All modalities share one normalization scheme (\texttt{normalization.py}), the
single source of truth for node identity. A normalized identifier is the
concatenation
\[
\texttt{normed\_id} \;=\; \text{normalized\_body}\;\Vert\;\text{``\_''}\;\Vert\;\text{md5(raw)[:6]},
\]
where the body lowercases the string, maps spaces to underscores, replaces any
character outside \texttt{[a-z0-9\_-]} with an underscore, collapses runs of
underscores, and preserves existing hyphens. Per-flavor front-ends specialize the
body (titles are capped at 193 characters; repository and package names split on
\texttt{/} and \texttt{.} first). The trailing 6-character MD5 disambiguates
strings that collide under the body rule---e.g.\ \texttt{"A+B"} and
\texttt{"A-B"}, both bodies \texttt{a\_b}. One asymmetry matters: arXiv ids are
\emph{version-stripped} to a canonical form \emph{before} hashing, so the hash is
taken over the canonical id (not the raw string); this guarantees a paper node
and every citation pointing at it produce the same identifier.

The consequential subtlety is a \textbf{raw-vs-canonical asymmetry}. The hashed
identifier is an \emph{internal} graph key, but the link targets that appear in
document \emph{text}---a wikilink to \texttt{[[Alan Turing]]}, or a Stack node
whose id embeds a raw file path---are the \emph{raw}, human-readable strings.
The runtime detectors accordingly match on the raw surface form (the Markdown
detector matches the raw title with spaces, not the hashed key), so node identity
and link surface form live in two spaces reconciled only through normalization.
If the normalization used at build time diverges by even one rule from that used
at comparison or merge time, targets stop matching node ids, every edge is
silently dropped, and the dataset comes out \emph{empty with no error raised}.
This is why the normalization module is centralized and treated as immutable
rather than forked per corpus.

\subsection{Pretokenization}
\label{app:datasets:pretok}

Node text is tokenized with \texttt{tiktoken}'s GPT-2 BPE (vocabulary 50{,}257,
stored as \texttt{uint16}). A pool of worker processes feeds a single writer over
a multiprocessing queue---a deliberate single-writer design that avoids the
interleaving deadlocks of a naive multi-writer scheme. Tokens are packed into
binary shards with a 2\,GB default target, each prefixed by a 1024-byte header
($256{\times}$\texttt{int32}) carrying a magic number, version, token count, and
element size. The tokenized graph merges each node's metadata with the shard
index, byte offset, and token length of its tokens. Consistent with the pipeline
invariant, link positions are not stored; the span realizing each edge is
re-detected at runtime (\Cref{app:linkdet}). This keeps the on-disk format
compact and modality-agnostic, at the cost of requiring the detectors to
reproduce exactly the link surface forms extraction assumed.

\subsection{Graph splitting}
\label{app:datasets:split}

Splitting (\texttt{split\_graph.py}) partitions nodes into \texttt{train},
\texttt{val\_random}, \texttt{val\_community}, \texttt{test\_random}, and
\texttt{test\_community} (default 2.5\% per held-out split). Each split is written
as a \emph{self-contained} directory with its own tokenized graph; crucially the
token shards are \emph{shared} by absolute-path reference and never copied, so the
five directories are cheap views over one set of shard bytes.

The two held-out kinds differ in structure. The \texttt{*\_random} splits are a
uniform sample of non-community nodes and carry no relational structure. The
\texttt{*\_community} splits are ``community packs'': held-out connected subgraphs
grown by \emph{bidirectional} BFS (following both out- and in-adjacency) from
seed nodes drawn in \emph{randomized} order, with each community constrained to
between 50 and 5{,}000 nodes; the pooled communities are shuffled and split 50/50
into \texttt{val\_community} and \texttt{test\_community}. Because a community
retains its intact intra-community link structure, validation loss on it reflects
genuine cross-document traversal rather than isolated-document perplexity, which
is the axis this paper cares about.

In every split, edges are filtered to endpoints \emph{within the same split};
any edge crossing a split boundary is severed. Membership is enforced at the node
level (each node belongs to exactly one split), so although all splits share the
same shard bytes, no held-out node's tokens are ever a training target and there
is no token-level leakage through the shared shards. Source stratification is
available as an option.

\subsection{Points a reviewer should weigh}
\label{app:datasets:attack}

The construction choices most open to challenge:

\begin{itemize}
  \item \textbf{Non-deterministic import resolution.} The unseeded tie-break
  means the Python graph is not reproducible run-to-run: two builds of the same
  repositories can yield different edge sets. A seeded or lexicographic tie-break
  would remove this.

  \item \textbf{Weak, sampling-only deduplication.} We rely on the upstream
  \texttt{the-stack-dedup} for code and perform \emph{no} additional dedup for
  Wikipedia or arXiv. The only in-house measurement is a whitespace-normalized
  prefix hash over a $\approx 1\%$ sample of Stack files, which is
  \emph{reported, not filtered}. There is no cross-split near-duplicate check, so
  a near-duplicate of a held-out document could remain in \texttt{train}; we do
  not quantify this, and robust decontamination against paraphrased
  near-duplicates \citep{yang2023rephrased} remains future work.

  \item \textbf{Selection bias from the two-link filter.} Discarding files with
  fewer than two intra-repo links removes the majority of Python files and
  over-represents package-heavy repositories; the highest-degree nodes tend to be
  \texttt{\_\_init\_\_.py} files and generated SDK code, which shapes what
  ``following an import'' means in aggregate.

  \item \textbf{Unmeasured Wikipedia resolution rate.} The wikilink resolution
  fraction is never logged, so we cannot state how many links became edges.
  Worse, a legacy fallback path invents link ids using a \emph{second}, different
  \texttt{md5[:6]} scheme (hashing a differently-cleaned string) that does not
  agree with the shared normalizer; any target routed through it is
  \emph{guaranteed} to dangle. The topology in \Cref{tab:datasets:scale} is thus
  a lower bound on the true reference density.

  \item \textbf{Incomplete import denylist and external arXiv coverage.} The
  denylist is not exhaustive: an unlisted third-party package (e.g.\
  \texttt{scipy}, \texttt{pydantic}) is treated as an intra-repo target and can
  produce spurious unresolved references. Separately, arXiv coverage climbs from
  $\approx 14.5\%$ to $\approx 66\%$ only via the external OpenAlex map
  \citep{priem2022openalex}, and therefore inherits that map's completeness and
  versioning.
\end{itemize}

The multi-source merge used for the diversity experiment (11 corpora unified by
normalized identifier with per-node source provenance and collision-safe id
remapping) reuses exactly the extraction, normalization, and pretokenization
machinery above; its treatment belongs with that experiment and is deferred to
\Cref{app:recipe} rather than repeated here.
